\chapter{Горизонт в плоских координатах и термодинамика испарения}

\lectureoverview{Метрика Шварцшильда, по-видимому, становится сингулярной на горизонте, хотя свободно падающий наблюдатель не должен заметить там ничего исключительного. Мы снимаем это противоречие, приводя пригоризонтную метрику к плоскому пространству-времени Минковского в гиперболических полярных координатах. Та же геометрия затем отделяет описание Алисы, падающей внутрь, от ускоренного внешнего описания Боба. Во второй половине вводятся энтропия и температура, выводятся температура Хокинга, обратно пропорциональная массе, и закон испарения \(M^3\), а в заключение ставится вопрос о том, чего сохранение информации требует от микроскопической теории гравитации.}

\section{Загадка при \texorpdfstring{\(r=2GM\)}{r=2GM}}

Начнем с очень большой черной дыры. Кривизна ее горизонта может быть настолько мала, что наблюдатель, падающий через него, в любой не слишком большой лабораторной области видит почти пустое плоское пространство-время. Он не натыкается ни на какую стену, и никакая локальная тревога не возвещает о пересечении. И все же метрика Шварцшильда, по-видимому, утверждает, что именно на этом радиусе происходит нечто сингулярное. Цель первой половины лекции — сделать эти два утверждения совместимыми.

Используем сигнатуру с одним минусом. В плоском пространстве-времени,
\begin{equation}
ds^2=-dt^2+dx^2+dy^2+dz^2,
\end{equation}
поэтому для времениподобной мировой линии \(ds^2<0\), а собственное время удовлетворяет \(d\tau^2=-ds^2\). Это просто общая смена знака по сравнению с соглашением, использованным в предыдущей лекции; ни одна наблюдаемая от этого выбора не зависит.

При \(c=1\) линейный элемент Шварцшильда имеет вид
\begin{equation}
\begin{aligned}
ds^2={}&-\left(1-\frac{2GM}{r}\right)dt^2
+\frac{dr^2}{1-\frac{2GM}{r}}
+r^2d\Omega_2^2,\\
d\Omega_2^2={}&d\theta^2+\sin^2\theta\,d\phi^2 .
\end{aligned}
\label{eq:l03-schwarzschild}
\end{equation}
Радиус Шварцшильда равен \(R_s=2GM\), или \(R_s=2GM/c^2\), если восстановить \(c\). Вдали, при \(r\gg R_s\), оба этих коэффициента стремятся к своим плоскопространственным значениям. Однако при \(r=R_s\) коэффициент при \(dt^2\) обращается в нуль, а коэффициент при \(dr^2\) расходится. Это выглядит драматично. Вопрос в том, драматична ли сама геометрия или только координаты.

\section{Два урока о координатах}

\subsection{Безобидное начало полярных координат}

На обычной евклидовой плоскости,
\begin{equation}
ds^2=dx^2+dy^2.
\end{equation}
введем полярные координаты
\begin{equation}
x=\rho\cos\theta,
\qquad
y=\rho\sin\theta .
\end{equation}
Малое смещение имеет радиальную длину \(d\rho\) и поперечную длину \(\rho\,d\theta\), следовательно,
\begin{equation}
ds^2=d\rho^2+\rho^2d\theta^2.
\label{eq:l03-polar}
\end{equation}
При \(\rho=0\) коэффициент при угловом члене обращается в нуль. Из-за этого центр плоскости не становится сингулярным. Это означает лишь, что там любому значению \(\theta\) соответствует одна и та же точка. Следовательно, необычный коэффициент метрики может указывать на необычную систему координат, а не на необычную физику.

\subsection{Гиперболические полярные координаты}

Теперь заменим евклидову плоскость \(1+1\)-мерным пространством-временем Минковского,
\begin{equation}
ds^2=-dT^2+dX^2.
\end{equation}
В правом клине \(X>|T|\) введем
\begin{equation}
X=\rho\cosh\omega,
\qquad
T=\rho\sinh\omega,
\qquad \rho>0 .
\label{eq:l03-rindler-transform}
\end{equation}
Обычное тождество \(\cos^2\theta+\sin^2\theta=1\) заменяется на
\begin{equation}
\cosh^2\omega-\sinh^2\omega=1.
\end{equation}
Тем самым линия с постоянным \(\rho\) есть гипербола
\begin{equation}
X^2-T^2=\rho^2.
\end{equation}
В отличие от периодического угла \(\theta\), быстротоподобная координата \(\omega\) изменяется от \(-\infty\) до \(+\infty\). При больших положительных \(\omega\) величины \(\cosh\omega\) и \(\sinh\omega\) становятся почти равными, и гипербола приближается к будущему светоподобному лучу \(X=T\).

Дифференцируя уравнение~\eqref{eq:l03-rindler-transform}, получаем
\begin{align}
dX&=\cosh\omega\,d\rho+\rho\sinh\omega\,d\omega,\\
dT&=\sinh\omega\,d\rho+\rho\cosh\omega\,d\omega .
\end{align}
Перекрестные члены сокращаются, и остается
\begin{equation}
ds^2=-\rho^2d\omega^2+d\rho^2.
\label{eq:l03-rindler-metric}
\end{equation}

\begin{figure}[htbp]
\centering
\includegraphics[width=0.72\textwidth]{lecture_03_figure_02.png}
\caption{Итоговая доска с гиперболическими координатами: \(X=\rho\cosh\omega\), \(T=\rho\sinh\omega\), построение единичной гиперболы и плоская метрика в виде \(ds^2=-\rho^2d\omega^2+d\rho^2\).}
\lectureframe{00:16:55.000}
\end{figure}

Уравнение~\eqref{eq:l03-rindler-metric} по-прежнему описывает совершенно плоское пространство-время. Его временной коэффициент обращается в нуль при \(\rho=0\), точно так же как угловой коэффициент в уравнении~\eqref{eq:l03-polar} обращается в нуль в начале координат. Координатная карта там вырождается; пространство-время — нет.

\section{Приближение к горизонту Шварцшильда}

Вернемся к радиально-временной части уравнения~\eqref{eq:l03-schwarzschild}:
\begin{align}
ds^2
&=-\frac{r-2GM}{r}\,dt^2
+\frac{r}{r-2GM}\,dr^2 .
\label{eq:l03-rewritten-schwarzschild}
\end{align}
Теперь рассмотрим только узкую окрестность \(r=2GM\). В множителях, которые меняются медленно, \(r\) можно заменить на \(2GM\). Множитель \(r-2GM\) необходимо сохранить, поскольку на горизонте он обращается в нуль и меняет знак. Поэтому
\begin{equation}
ds^2\approx
-\frac{r-2GM}{2GM}\,dt^2
+\frac{2GM}{r-2GM}\,dr^2 .
\label{eq:l03-near-horizon-r}
\end{equation}

Выберем новую радиальную координату \(\rho\) так, чтобы второй член был просто \(d\rho^2\):
\begin{equation}
d\rho^2=\frac{2GM}{r-2GM}\,dr^2.
\end{equation}
Вне горизонта, взяв положительный квадратный корень и проинтегрировав от \(r=2GM\), получаем
\begin{align}
d\rho&=\sqrt{\frac{2GM}{r-2GM}}\,dr,\\
\rho&=2\sqrt{2GM}\sqrt{r-2GM},\\
\rho^2&=8GM(r-2GM),
\qquad
r-2GM=\frac{\rho^2}{8GM}.
\label{eq:l03-rho-r}
\end{align}
Подстановка в уравнение~\eqref{eq:l03-near-horizon-r} дает
\begin{equation}
ds^2\approx
-\frac{\rho^2}{16G^2M^2}\,dt^2+d\rho^2.
\end{equation}
Наконец, переопределим время Шварцшильда,
\begin{equation}
\omega=\frac{t}{4GM},
\qquad
d\omega^2=\frac{dt^2}{16G^2M^2},
\label{eq:l03-omega}
\end{equation}
и получим
\begin{equation}
\boxed{ds^2\approx-\rho^2d\omega^2+d\rho^2.}
\label{eq:l03-near-horizon-rindler}
\end{equation}

\begin{figure}[htbp]
\centering
\includegraphics[width=0.70\textwidth]{lecture_03_figure_03.png}
\caption{Доска после перескалирования времени: пригоризонтная метрика Шварцшильда превратилась в ту же гиперболически-полярную метрику, что и у плоского пространства-времени Минковского.}
\lectureframe{00:31:35.000}
\end{figure}

По мере того как рассматриваемая окрестность сжимается к горизонту, приближение становится сколь угодно точным. Это доказывает искомое локальное утверждение: для достаточно большой дыры горизонт — это обычное пространство-время, записанное в необычной координатной карте. Кривизна для шварцшильдовской дыры конечной массы не равна буквально нулю, но на горизонте она конечна и может быть сделана там сколь угодно малой, если увеличивать \(M\).\editorialnote{Локальный масштаб кривизны имеет порядок \(R_s^{-2}\). Метрика Риндлера является главным членом локального разложения вблизи горизонта, а не утверждением, что вся геометрия Шварцшильда глобально плоская.}

\begin{classroomqa}
\audiencequestion{Для \(r<2GM\) уравнение~\eqref{eq:l03-rho-r} делает внешнюю координату \(\rho\) мнимой. Указывает ли это на физическую проблему?\lecturetimestamp{00:32:30}}
\lecturerresponse{Это воспроизводит тот же самый сбой координат в другой форме. Данная карта была построена для внешнего клина; для продолжения через горизонт требуется другая карта. Сбой именно этой \(\rho\) не является физическим событием на горизонте.}
\end{classroomqa}

\begin{classroomqa}
\audiencequestion{При переписывании вблизи горизонта как приближенная строка получается из предыдущей точной строки?\lecturetimestamp{00:39:21}}
\lecturerresponse{Только те вхождения \(r\), которые меняются медленно, были приравнены \(2GM\). Малая разность \(r-2GM\) оставлена без изменения. В этом и состоит приближение, и оно становится все точнее по мере того, как \(r\) приближается к \(2GM\).}
\end{classroomqa}

\section{Что означают координаты}

Преобразование устанавливает
\begin{equation}
\rho=0\quad\Longleftrightarrow\quad r=2GM.
\end{equation}
Внешняя область \(r>2GM\) — это правый клин Риндлера \(X>|T|\). Гиперболы с постоянным \(\rho\) представляют наблюдателей, остающихся на фиксированном радиусе Шварцшильда. Исходящий световой луч, испущенный где угодно в этом клине, пересекает каждую мировую линию с большим \(\rho\), так что внешний наблюдатель может передавать сигналы наружу. Точка за будущей светоподобной границей не может послать световой сигнал обратно в правый клин. Эта светоподобная граница и есть горизонт.

\begin{classroomqa}
\audiencequestion{Всегда ли показанная здесь гиперболическая координатная область описывает область вне горизонта?\lecturetimestamp{00:35:22}}
\lecturerresponse{Да. В этой конструкции \(r>2GM\) дает \(\rho^2>0\) и помещает событие в правый внешний клин. Светоподобная граница \(\rho=0\) — это горизонт; данные координаты не продолжаются за него как вещественные внешние координаты.}
\end{classroomqa}

\begin{classroomqa}
\audiencequestion{Представляет ли гипербола с постоянным \(\rho\) время или расстояние?\lecturetimestamp{00:37:21}}
\lecturerresponse{Она задает одно фиксированное собственное расстояние до горизонта в пригоризонтном приближении и, следовательно, один фиксированный радиус Шварцшильда. Движение вдоль гиперболы есть течение времени для находящегося там наблюдателя.}
\end{classroomqa}

Вдоль траектории с фиксированным \(\rho\) уравнение~\eqref{eq:l03-rindler-metric} дает
\begin{equation}
d\tau=\rho\,d\omega.
\label{eq:l03-proper-time-rindler}
\end{equation}
Следовательно, \(\omega\) само по себе не является собственным временем наблюдателя. Это угловое время, общее для семейства гипербол; каждый наблюдатель умножает его на свое \(\rho\), точно так же как обычный угол умножается на радиус окружности, чтобы дать длину дуги.

\begin{classroomqa}
\audiencequestion{Является ли \(\omega\) собственным временем ускоренного наблюдателя?\lecturetimestamp{00:51:45}}
\lecturerresponse{Почти. При фиксированном \(\rho\) приращение собственного времени равно \(\rho\,d\omega\). Множитель \(\rho\) существен.}
\end{classroomqa}

\begin{figure}[htbp]
\centering
\includegraphics[width=0.72\textwidth]{lecture_03_figure_04.png}
\caption{Пригоризонтная причинная диаграмма на доске: пересекающиеся светоподобные границы, гиперболы постоянного радиуса во внешнем клине и координатная связь между \(\rho\) и \(r-2GM\).}
\lectureframe{00:34:35.000}
\end{figure}

\begin{figure}[htbp]
\centering
\begin{tikzpicture}[scale=1.04,>=Latex]
  \draw[->] (-0.3,0) -- (4.6,0) node[right] {\(X\)};
  \draw[->] (0,-3.2) -- (0,3.2) node[above] {\(T\)};
  \draw[very thick,blue!55!black] (0,0) -- (3,3);
  \draw[very thick,blue!55!black] (0,0) -- (3,-3);
  \node[blue!55!black,rotate=45] at (0.72,1.02) {future horizon};
  \draw[thick,domain=-1.45:1.45,samples=80] plot ({0.75*cosh(\x)},{0.75*sinh(\x)});
  \draw[thick,domain=-1.25:1.25,samples=80] plot ({1.25*cosh(\x)},{1.25*sinh(\x)});
  \draw[very thick,domain=-1.05:1.05,samples=80] plot ({2.0*cosh(\x)},{2.0*sinh(\x)});
  \node at (3.64,2.18) {Bob: \(\rho=\mathrm{const.}\)};
  \draw[very thick,red!65!black,->] (1.2,-2.9) -- (1.2,2.9);
  \node[red!65!black] at (1.58,-1.75) {Alice};
  \fill[red!65!black] (1.2,1.2) circle (1.7pt);
  \draw[dashed,->] (1.2,0) -- (2.27,1.07);
  \node[align=left] at (3.18,0.80) {\small pre-crossing\\[-0.2em]\small light signal};
\end{tikzpicture}
\caption{Набранное дополнение к доске. Боб движется по ускоренной гиперболе с постоянным \(\rho\); Алиса движется по локально инерциальной линии через будущий горизонт. Светоподобная граница ограничивает причинный доступ Боба, не становясь локальным препятствием для Алисы.}
\end{figure}

\section{Алиса падает, а Боб удерживается на месте}

Пусть Боб удерживается на фиксированном расстоянии от горизонта, а Алиса отпускает его руку и падает. Мировая линия Боба — гипербола и, следовательно, ускоренная. Локальная мировая линия Алисы приблизительно прямая. Она пересекает светоподобную границу за конечное собственное время и не встречает никакого локального разрыва.

Причинные описания несимметричны. После того как Алиса пересекает горизонт, свет от неё уже не может достичь Боба, тогда как свет, испущенный Бобом, ещё некоторое время может достигать Алисы. Боб получает сигналы, испущенные Алисой на всё более ранних этапах её конечного падения. Интервалы между приходами этих сигналов становятся всё больше, их частота смещается в красную сторону, а интенсивность ослабевает. В шварцшильдовом времени он так и не видит самого пересечения; кажется, что Алиса лишь асимптотически приближается к горизонту и распластывается по нему. Дело не в том, что у неё расплющивается нос, тогда как затылок остаётся обычным: сигналы от всех частей сжимаются в один и тот же предельный слой.

\begin{classroomqa}
\audiencequestion{Может ли Алиса видеть Боба после пересечения, потому что фотоны Боба летят внутрь к ней?\lecturetimestamp{00:42:48}}
\lecturerresponse{Да. В этом и состоит односторонность причинной связи: входящие сигналы Боба могут проникать внутрь, тогда как более поздние исходящие сигналы Алисы уже не могут выбраться обратно к Бобу.}
\end{classroomqa}

\begin{classroomqa}
\audiencequestion{Если ноги Алисы пересекают горизонт первыми, тогда как её глаза остаются снаружи, теряет ли она временно из виду собственные ноги?\lecturetimestamp{00:43:03}}
\lecturerresponse{При обычном падении никакого зрительного сбоя не происходит. Свету всегда нужно время, чтобы пройти от ног к глазам; к тому времени, когда приходит соответствующий сигнал, горизонт уже пересекли и её глаза. Лишь грубая попытка удерживать её голову снаружи, пока ноги продолжают уходить внутрь, создала бы иную и разрушительную ситуацию.}
\end{classroomqa}

Удерживаться на месте — значит сделать неинерциальный выбор. Бобу нужна ракета, точно так же как человеку, стоящему на полу, нужна направленная вверх сила со стороны пола, чтобы избежать свободного падения. Чем ближе Боб зависает у горизонта, тем больше требуется собственное ускорение; в локальном риндлеровском приближении его величина равна \(a=1/\rho\) в единицах \(c=1\).\editorialnote{При восстановлении единиц для локального риндлеровского наблюдателя получаем \(a=c^2/\rho\). У строго статического шварцшильдовского наблюдателя собственное ускорение расходится при приближении к горизонту.}

\begin{classroomqa}
\audiencequestion{Ускоряемся ли мы, стоя здесь, в аудитории?\lecturetimestamp{00:45:40}}
\lecturerresponse{Да, относительно свободно падающей системы отсчёта. Пол нас поддерживает; без него или без реактивного ранца мы бы падали.}
\end{classroomqa}

\begin{classroomqa}
\audiencequestion{Увидит ли Боб в конце концов Алису в пределах одной планковской длины от горизонта, и что происходит после этого?\lecturetimestamp{00:49:24}}
\lecturerresponse{Именно здесь классического геометрического описания уже недостаточно. В планковскую длину входит \(\hbar\), так что для ответа на вопрос, что именно Боб там может различить, требуется квантовая механика.}
\end{classroomqa}

\begin{classroomqa}
\audiencequestion{Кривые постоянного радиуса как будто встречаются с горизонтом в координатном углу, похожем на прямой. Можно ли в такой точке брать производную?\lecturetimestamp{00:51:12}}
\lecturerresponse{В этой точке нечто сингулярное происходит с данной координатной картой, но не с локальной физикой. Карта, покрывающая другие клинья, убирает этот искусственный угол. Для текущего обсуждения внешней области риндлеровская картина уже содержит все нужные причинные факты.}
\end{classroomqa}

Поэтому бесполезно, не указав наблюдателя, спрашивать, пересекает ли Алиса горизонт \emph{на самом деле}. По часам Алисы горизонт достигается за конечное собственное время. В последовательности принимаемых Бобом сигналов событие пересечения никогда не появляется. Физика спрашивает, что измеряет каждый наблюдатель и как сопоставить их записи.

\begin{classroomqa}
\audiencequestion{Верно ли точное утверждение, что Боб может видеть лишь конечную часть будущего Алисы, хотя его собственное время наблюдения неограниченно?\lecturetimestamp{00:53:33}}
\lecturerresponse{Да. Алиса приписывает видимой части своего падения конечное собственное время. В идеализированной вечной геометрии Боб бесконечно долго получает её всё более запаздывающие и всё более красносмещённые сигналы, но не видит, как она продолжает стареть после пересечения горизонта.}
\end{classroomqa}

\section{Одних горизонтов ускорения уже достаточно}

В этой причинной структуре ничто не требует наличия чёрной дыры. Ракета, которая равномерно ускоряется вечно, движется по гиперболе Риндлера в плоском пространстве-времени и имеет светоподобную границу, по ту сторону которой никакой сигнал не может её догнать. Для такого наблюдателя существует горизонт ускорения, хотя инерциальный наблюдатель считает ту же область обычным пространством Минковского.

\begin{classroomqa}
\audiencequestion{Это явление специфично для чёрных дыр или оно возникает при любом горизонте?\lecturetimestamp{00:55:25}}
\lecturerresponse{То же самое явление возникает и у равномерно ускоренного наблюдателя. У такого наблюдателя есть горизонт ускорения, и он не может видеть события за ним. Локально горизонт чёрной дыры физически не отличается от этого случая.}
\end{classroomqa}

Так проявляется принцип эквивалентности: сначала разберитесь с равномерным ускорением, а затем распознайте ту же локальную структуру в гравитационном поле. Оглядываясь назад, можно было предвидеть гравитационные горизонты, исходя из горизонтов ускорения, ещё до того, как чёрные дыры были признаны. Сам Эйнштейн не принимал возможность физического образования чёрных дыр, хотя именно он дал теорию, из которой следует их геометрия. Для свободно падающего наблюдателя горизонт скучен, а для наблюдателя, который настаивает на том, чтобы оставаться снаружи, он имеет последствия. Это напряжение между корректными описаниями, а не парадокс — пока ещё.

\section{Возвращаясь к теплоте и энтропии}

После перерыва наденем термодинамическую шляпу. В предыдущем построении чёрную дыру наполняли, как ванну, по одной капле, с той лишь разницей, что каждая капля была фотоном, несущим примерно один бит. Каждый бит увеличивал площадь горизонта на величину порядка планковской площади. Точная энтропия равна
\begin{equation}
S_{\mathrm{BH}}=\frac{k_B A}{4\ell_P^2}.
\end{equation}
В лекции шутят, что множителя \(1/4\) можно было бы избежать, если бы Планк определил свою длину иначе. Серьёзный смысл в том, что размерный анализ фиксирует масштаб, но не численный коэффициент; коэффициент задаётся вычислением Хокинга.\editorialnote{Планковская длина стандартно определяется как \(\ell_P=\sqrt{G\hbar/c^3}\). Множитель \(1/4\) имеет физический смысл, как только зафиксированы это соглашение и обычная термодинамическая энтропия; замечание о его переопределении — аудиторная шутка о соглашениях.}

Температура кажется чем-то элементарным, потому что её напрямую регистрируют пальцы и термометры. Энтропия же часто кажется более загадочной. Лекция переворачивает эту иерархию: энтропия и энергия рассматриваются как основные переменные, а температура выводится из того, насколько должна измениться энергия при изменении энтропии:
\begin{equation}
dE=T\,dS.
\end{equation}
В используемых в лекции единицах, где \(\Delta S=1\), минимальная энергия, связанная с добавлением или стиранием одного бита, обозначается как \(T\). Небрежное стирание --- логическим эквивалентом кувалды --- может рассеять сколь угодно большее количество энергии. Это определение относится к наименьшей возможной диссипации, когда отсутствующая информация передаётся микроскопическим степеням свободы окружающей среды. В обычных единицах один двоичный бит несёт \(\Delta S=k_B\ln2\), а обратимое стирание рассеивает по меньшей мере \(k_BT\ln2\).\editorialnote{Это и есть нормировка Ландауэра, скрытая тем, что в лекции и энтропия, и температура измеряются в естественных единицах, основанных на энергии.}

\subsection{Однофотонная оценка}

Выберем фотон, длина волны которого сравнима с радиусом Шварцшильда:
\begin{equation}
\lambda\sim R_s=\frac{2GM}{c^2}.
\end{equation}
Его энергия с точностью до соглашения относительно \(2\pi\), различающего \(h\) и \(\hbar\), равна
\begin{equation}
\Delta E\sim \frac{\hbar c}{\lambda}.
\end{equation}
Это и есть энергетическая цена добавленного бита, так что
\begin{equation}
k_BT_H\sim\frac{\hbar c}{R_s}
\sim\frac{\hbar c^3}{GM}.
\label{eq:l03-hawking-scale}
\end{equation}
Эвристический аргумент не может определить множители вроде \(2\) и \(\pi\). Точный результат Хокинга таков:
\begin{equation}
\boxed{k_BT_H=\frac{\hbar c^3}{8\pi GM}.}
\label{eq:l03-hawking-exact}
\end{equation}

\begin{figure}[htbp]
\centering
\includegraphics[width=0.72\textwidth]{lecture_03_figure_05.png}
\caption{Итоговая запись на доске для температуры: \(R_s=2GM/c^2\), масштаб одного бита \(\hbar c/R_s\) и точная хокинговская нормировка \(8\pi\).}
\lectureframe{01:07:48.000}
\end{figure}

Сразу видны два факта. Во-первых, температура обращается в нуль вместе с \(\hbar\), так что это квантовый эффект. Во-вторых,
\begin{equation}
T_H\propto\frac{1}{M}.
\end{equation}
Чёрная дыра большой массы холодна. Для чёрной дыры солнечной массы
\begin{equation}
T_H\approx6\times10^{-8}\ {\rm K},
\end{equation}
что имеет тот же порядок, что и лекционная оценка \(10^{-8}\,\mathrm{K}\), и намного ниже температуры космического микроволнового фона, составляющей около \(2.7\,\mathrm{K}\).

\subsection{Отрицательная теплоёмкость}

Если поместить такую чёрную дыру в более тёплую среду, она поглощает энергию, её масса растёт, и из уравнения~\eqref{eq:l03-hawking-exact} следует, что она становится холоднее. После этого она поглощает ещё охотнее. Если поместить её в более холодную среду, начинается обратный неустойчивый процесс: она излучает, теряет массу, становится горячее и излучает быстрее. Формально,
\begin{equation}
C_{\mathrm{BH}}
=\frac{dE}{dT_H}
=c^2\frac{dM}{dT_H}
<0.
\end{equation}

\begin{classroomqa}
\audiencequestion{Это и имеется в виду под отрицательной теплоёмкостью?\lecturetimestamp{01:14:08}}
\lecturerresponse{Именно. Самогравитирующая система может нагреваться, теряя энергию. Со звездой происходит нечто аналогичное: по мере излучения уменьшение давления позволяет ей сжиматься, а гравитационное сжатие повышает её температуру.}
\end{classroomqa}

Изолированная чёрная дыра от этого не становится нефизической. В современной Вселенной чёрная дыра звёздной массы холоднее своего окружения и преимущественно поглощает. По мере расширения Вселенной и охлаждения фона это соотношение в конечном счёте может измениться на противоположное, и испарение может начать доминировать. Простое равновесие с эффективно бесконечным термостатом неустойчиво.\editorialnote{Граничные условия важны. Чёрную дыру можно стабилизировать в конечной отражающей полости или, для достаточно больших дыр, в пространстве-времени анти-де Ситтера; описанная здесь неустойчивость относится к асимптотически плоскому случаю с обычным термостатом.}

\section{Закон испарения}

Для вывода по порядку величины представьте, что у вас нет под рукой справочника, и выберите единицы, в которых вспоминаемые уравнения имеют самый короткий вид:
\begin{equation}
c=\hbar=G=k_B=1.
\end{equation}
Это планковские единицы. Характерные масштабы, приведённые в лекции, таковы:
\begin{equation}
\ell_P\sim10^{-35}\ {\rm m},
\qquad
t_P\sim10^{-43}\ {\rm s},
\qquad
m_P\sim10^{-8}\ {\rm kg}.
\end{equation}
Планковская масса по порядку величины сопоставима с видимой пылинкой, тогда как планковские длина и время чрезвычайно малы. Поскольку \(c=1\), \(t_P\) — это просто время пролёта света через \(\ell_P\).

Рассматривайте дыру как абсолютно чёрное тело, то есть как тепловой излучатель, а не как объект, который обязан оставаться совершенно тёмным. Она может испускать электромагнитное излучение и, в принципе, другие доступные типы частиц, включая гравитоны. Излучаемая мощность пропорциональна площади излучающей поверхности, потому что излучение уходит через поверхность. Масштабирование Стефана--Больцмана имеет вид
\begin{equation}
L=-\frac{dE}{dt}\sim A T^4.
\end{equation}
В планковских единицах \(E=M\), \(R_s\sim M\), \(A\sim R_s^2\sim M^2\) и \(T\sim R_s^{-1}\sim M^{-1}\). Поэтому
\begin{align}
-\frac{dM}{dt}
&\sim R_s^2\left(\frac{1}{R_s}\right)^4
\sim\frac{1}{R_s^2}
\sim\frac{1}{M^2},\\
-M^2\frac{dM}{dt}&\sim1,\\
-\frac{d(M^3)}{dt}&\sim1.
\label{eq:l03-evap-scaling}
\end{align}
Если пренебречь коэффициентами порядка единицы, время жизни равно
\begin{equation}
\boxed{t_{\mathrm{evap}}\sim M^3.}
\label{eq:l03-lifetime}
\end{equation}

\begin{figure}[htbp]
\centering
\includegraphics[width=0.72\textwidth]{lecture_03_figure_06.png}
\caption{Вывод на доске для беспрепятственного испарения: светимость даёт \(-dM/dt\sim1/M^2\), отсюда равномерное убывание \(M^3\) и время жизни, пропорциональное кубу начальной массы.}
\lectureframe{01:23:20.000}
\end{figure}

В схематическом рассмотрении по Стефану--Больцману опущены серотельные факторы и число видов частиц. Для нейтральной невращающейся дыры в простейшем приближении излучения безмассовых частиц восстановление констант даёт
\begin{equation}
t_{\mathrm{evap}}
=\frac{5120\pi G^2M_0^3}{\hbar c^4}.
\label{eq:l03-exact-lifetime}
\end{equation}
Этот коэффициент — редакторское дополнение; в лекции выводится только надёжная зависимость вида \(M_0^3\).

\subsection{Масштабы солнечной массы, горы и килограмма}

Солнце имеет массу \(M_\odot\sim10^{30}\,\mathrm{kg}\sim10^{38}m_P\). Возведение в куб и умножение на одно планковское время даёт намеренно грубую цепочку оценок из лекции
\begin{equation}
t_{\mathrm{evap}}\sim10^{114}t_P
\sim10^{71}\ {\rm s}
\sim10^{64}\ {\rm yr}.
\end{equation}
Обе оценки фантастически превышают возраст Вселенной, составляющий примерно \(10^{10}\) лет. Восстановление коэффициента в уравнении~\eqref{eq:l03-exact-lifetime} увеличивает стандартную оценку для солнечной массы примерно до \(2\times10^{67}\) лет.\editorialnote{Это различие ожидаемо: в вычислении на доске отброшены все численные константы, а планковские масштабы округлены до степеней десяти. Его цель — зависимость как куб массы, а не прецизионные астрономические оценки времени жизни.}

Затем в аудитории рассматривают чёрную дыру \emph{горной массы}. Останавливаются на грубом аудиторном значении \(M\sim10^{10}\,\mathrm{kg}\), откуда \(M\sim10^{18}m_P\), и получают
\begin{equation}
t_{\mathrm{evap}}\sim10^{54}t_P
\sim10^{11}\ {\rm s}
\sim10^4\ {\rm yr}.
\end{equation}
Здесь сохранена арифметика, действительно проделанная на доске. Словесная геометрическая оценка, использованная для обоснования этой массы---примерно \(100\,\mathrm{km}^3\) каменной породы---, в действительности дала бы гораздо большую массу, а точный коэффициент времени жизни тоже существенно изменил бы результат.\editorialnote{При обычной плотности горных пород \(100\,\mathrm{km}^3\) имеет порядок \(10^{14}\,\mathrm{kg}\), а не \(10^{10}\,\mathrm{kg}\). Для \(10^{10}\,\mathrm{kg}\) уравнение~\eqref{eq:l03-exact-lifetime} даёт время жизни порядка \(10^6\) лет. Эти поправки не меняют урок лекции о том, что время жизни масштабируется как куб массы.}

\begin{classroomqa}
\audiencequestion{Будет ли чёрная дыра массой всего в килограмм чрезвычайно маленькой?\lecturetimestamp{01:29:14}}
\lecturerresponse{Да. Поскольку радиус пропорционален массе, а солнечной массе соответствует несколько километров, радиус Шварцшильда для одного килограмма составляет около \(10^{-27}\,\mathrm{m}\). Такая дыра также была бы чрезвычайно горячей и короткоживущей.}
\end{classroomqa}

\begin{classroomqa}
\audiencequestion{Предполагают ли эти оценки испарения нулевую температуру окружающей среды?\lecturetimestamp{01:29:53}}
\lecturerresponse{Они применимы, когда окружение холоднее дыры, так что суммарный поток энергии направлен наружу. Более горячая среда, напротив, подпитывает чёрную дыру. Масса, при которой температура Хокинга совпадает с нынешним фоном, оставляется в качестве упражнения.}
\end{classroomqa}

Для ориентира, если положить \(T_H=2.7\,\mathrm{K}\), то \(M\) имеет порядок \(10^{22}\,\mathrm{kg}\), а радиус Шварцшильда — порядок \(10^{-4}\,\mathrm{m}\).\editorialnote{Здесь вычислено упражнение, предложенное в лекции; в устный вывод это не входило.}

\section{Что должно быть под метрикой?}

\begin{classroomqa}
\audiencequestion{Какой механизм заставляет фотоны выходить из чёрной дыры?\lecturetimestamp{01:30:53}}
\lecturerresponse{Прежде чем давать микроскопический механизм излучения, нужно определить степени свободы, ответственные за энтропию. Термодинамический аргумент накладывает ограничения на то, что должно существовать, ещё не выписывая эти переменные явно.}
\end{classroomqa}

Здесь уместна аналогия с механикой жидкости. Гидродинамика описывает гладкие поля плотности и скорости, не отслеживая каждую молекулу. Тем не менее энтропия показывает, что одному и тому же гладкому состоянию соответствуют многие скрытые микроскопические конфигурации. Точно так же классическая общая теория относительности описывает гладкую метрику, но не раскрывает многочисленных микроскопических степеней свободы, учитываемых энтропией чёрной дыры. Закон площади и хокинговское излучение служат указаниями на то, что теория Эйнштейна является эффективным макроскопическим описанием.

\begin{classroomqa}
\audiencequestion{Если чёрная дыра излучает энергию, означает ли это то же самое, что она излучает информацию, так что информация, поглощённая при её росте, возвращается при испарении?\lecturetimestamp{01:33:27}}
\lecturerresponse{Таково требование: если полная эволюция сохраняет информацию, исходящее излучение в конечном счёте должно различать возможные состояния, из которых образовалась дыра. Это ещё нужно согласовать с причинной картиной, в которой вещество пересекает горизонт.}
\end{classroomqa}

Сам по себе поток энергии недостаточен, чтобы задать информацию. Совершенно тепловая матрица плотности может переносить энергию, забывая при этом подробное начальное состояние. В унитарном описании требуемая информация закодирована в тонких корреляциях между испущенными квантами.\editorialnote{Это различие явно формулирует современную форму проблемы информации. Оно не меняет, а заостряет вопрос, поставленный в лекции.}

\section{Сохранение информации как структурный принцип}

\begin{classroomqa}
\audiencequestion{Действительно ли сохранение информации является фундаментальным принципом?\lecturetimestamp{01:34:13}}
\lecturerresponse{Оно относится к числу наиболее фундаментальных. Если бы различные начальные состояния могли эволюционировать в совершенно одно и то же конечное состояние, законы стирали бы своё собственное прошлое. Обычные законы сохранения опираются не только на симметрии, но и на обратимую гамильтоновскую либо унитарную эволюцию.}
\end{classroomqa}

Вообразите, как это делается в лекции, вымышленный фундаментальный закон, который просто заставляет любой скользящий ластик останавливаться. Все разные начальные скорости приводят к одному и тому же состоянию. В конечном состоянии ластик уже не различает свои предыстории: информация уничтожена. Исчезают также импульс и механическая энергия. Настоящее трение на микроскопическом уровне так не действует; импульс и энергия переходят в стол, воздух и тепловые степени свободы.

Даже совершенно однородная бесконечная столешница обладала бы симметрией переносов, и всё же фиктивный закон затухания мог бы заставить ластик остановиться. Следовательно, одной лишь симметрии переносов недостаточно, чтобы вывести из такого негамильтонова правила закон сохранения импульса. Теорема Нётер предполагает принцип действия и связанную с ним обратимую динамику. На языке, используемом здесь, именно симметрия вместе с динамикой, сохраняющей информацию, даёт закон сохранения.

\begin{classroomqa}
\audiencequestion{Является ли сохранение информации более фундаментальным, чем причинность?\lecturetimestamp{01:37:00}}
\lecturerresponse{Это разные принципы сопоставимого уровня фундаментальности. Сначала нужно точно сформулировать причинность---например, что сигналы не могут распространяться назад за пределами допустимой причинной структуры---но ни один из этих принципов не может просто заменить другой.}
\end{classroomqa}

\begin{classroomqa}
\audiencequestion{Является ли обсуждаемое в гамильтоновой механике сохранение площади фазового пространства классической версией этой идеи?\lecturetimestamp{01:37:38}}
\lecturerresponse{Да. Теорема Лиувилля утверждает, что гамильтонов поток сохраняет объём фазового пространства. Облако возможных начальных состояний не может сжаться в меньший объём и стереть различия. В квантовой механике соответствующее утверждение — унитарность: сохраняются скалярные произведения, а значит, и различия между ортогональными состояниями.}
\end{classroomqa}

Классически,
\begin{equation}
\frac{d}{dt}\,\mathrm{Vol}_{\Gamma}=0,
\end{equation}
для области фазового пространства, переносимой гамильтоновым потоком. В квантовой механике,
\begin{equation}
\begin{aligned}
\ket{\psi(t)}&=U(t)\ket{\psi(0)},\\
U^\dagger U&=I,\\
\braket{\psi(t)}{\phi(t)}
&=\braket{\psi(0)}{\phi(0)}.
\end{aligned}
\end{equation}
Это точные формы утверждения о том, что микроскопические различия сохраняются при эволюции во времени.

\section{От чего зависит энтропия?}

В заключительном обмене репликами обсуждение возвращается к информационно-теоретической энтропии. Если вероятности микроскопических альтернатив обозначены как \(p_i\),
\begin{equation}
S=-k_B\sum_i p_i\ln p_i
\label{eq:l03-shannon}
\end{equation}
в термодинамических единицах, а без \(k_B\) — для безразмерной энтропии Шеннона.

\begin{figure}[htbp]
\centering
\includegraphics[width=0.72\textwidth]{lecture_03_figure_07.png}
\caption{Заключительная формула на доске \(S=-\sum_i p_i\log p_i\), записанная при проведении различия между распределением вероятностей и единственным значением энтропии, вычисленным из него.}
\lectureframe{01:40:32.000}
\end{figure}

\begin{classroomqa}
\audiencequestion{Определяет ли максимизация энтропии распределение вероятностей, а не вероятности — энтропию?\lecturetimestamp{01:38:33}}
\lecturerresponse{Распределение вероятностей задаёт одно значение энтропии по уравнению~\eqref{eq:l03-shannon}; по одному числу нельзя восстановить все \(p_i\). Максимизация энтропии может выделить распределение только после того, как заданы также нормировка и другие ограничения, например фиксированная средняя энергия.}
\end{classroomqa}

Этим различием лекция завершается у правой границы. Гладкая геометрия привела к горизонту; горизонт привёл к термодинамике; термодинамика выявила скрытые степени свободы; и сохранение информации теперь требует, чтобы эти степени свободы эволюционировали, не стирая различий, которые были у того, что вошло в дыру. Механизм пока ещё не указан. Вместо этого геометрия и учёт прояснили, что должна объяснить любая успешная микроскопическая теория.
